% !TEX root = task-2-5-4.tex
% Editable task bank. Append new BankTask blocks at the end of the document.
\documentclass[12pt,a4paper]{article}
\usepackage{physics-textbook}
\usepackage{graphicx,needspace}
\usetikzlibrary{decorations.pathmorphing}
\newif\ifShowSolutions
\ifdefined\TutorEdition
  \ShowSolutionstrue
\else
  \ShowSolutionsfalse % Change to \ShowSolutionstrue to display solutions.
\fi
\newcounter{banktask}
\newsavebox{\BankTaskBox}
\newcommand{\BankTaskContent}[3]{%
  \noindent\textbf{\thebanktask.}\quad #2\par
  \ifShowSolutions
    \Needspace{50mm}% Keep the solution heading with its opening text or diagram.
    #3\par
    \ifdefined\TutorEdition
      \ifcsname TutorSource#1\endcsname
        \nopagebreak\smallskip{\small\csname TutorSource#1\endcsname\par}%
      \fi
    \fi
  \fi
}
\newcommand{\BankTask}[3]{%
  \refstepcounter{banktask}\par\medskip
  \sbox{\BankTaskBox}{\begin{minipage}{\linewidth}%
    \BankTaskContent{#1}{#2}{#3}\end{minipage}}%
  \ifdim\dimexpr\ht\BankTaskBox+\dp\BankTaskBox\relax>\dimexpr\textheight-20mm\relax
    \Needspace{35mm}\BankTaskContent{#1}{#2}{#3}%
  \else
    \Needspace{\dimexpr\ht\BankTaskBox+\dp\BankTaskBox\relax}%
    \noindent\usebox{\BankTaskBox}\par
  \fi
  \medskip
}
\begin{document}
\BookSubsection{2.5.4}{Уравнение Менделеева — Клапейрона}


% Source sdamgia-813: https://phys-ege.sdamgia.ru/problem?id=813
\BankTask{sdamgia-813}{
При неизменной концентрации молекул абсолютная температура идеального газа была увеличена в $4$ раза. Во сколько раз изменилось давление газа?
}{
\smallskip\noindent\textbf{Решение}\par
Давление идеального газа пропорционально произведению концентрации молекул газа и абсолютной температуры: 

$p=nkT.$

При неизменной концентрации молекул и увеличении абсолютной температуры в $4$ раза давление газа также увеличится в $4$ раза.

Ответ: $4$.
}

% Source sdamgia-3105: https://phys-ege.sdamgia.ru/problem?id=3105
\BankTask{sdamgia-3105}{
В сосуде под поршнем находится $3$ моля гелия. Что произойдет с давлением газа на стенки сосуда, температурой и объемом газа при его изотермическом расширении?

\TaskMatchingColumns{Физические величины}{А)  Давление газа\par Б)  Температура газа\par В)  Объем газа}{Изменение величины}{$1$.  Увеличивается\par $2$.  Уменьшается\par $3$.  Не изменится}

\begin{center}\begin{tabular}{|c|c|c|}\hline
$A$ & Б & В\\\hline
\strut & \strut & \strut\\\hline\end{tabular}\end{center}
}{
\smallskip\noindent\textbf{Решение}\par
Изотермическим называется процесс при постоянной температуре. Следовательно, температура газа останется неизменной (Б — $3$). Поскольку газ расширяется в сосуде, его объем увеличивается (В — $1$). Гелий является инертным газом, поэтому его можно считать идеальным. При изотермическом процессе, согласно закону Бойля — Мариотта, величина $pV$ остается постоянной. Таким образом, заключаем, что при изотермическом расширении гелия в сосуде его давление уменьшается (А — $2$).

Ответ: $231$.
}

% Source sdamgia-19725: https://phys-ege.sdamgia.ru/problem?id=19725
\BankTask{sdamgia-19725}{
\begin{center}\begin{tikzpicture}[task graph]
\draw[task grid,step=.6] (0,0) grid (3.6,3.6);\draw[task axis,->] (-.2,0)--(4,0) node[right] {$\rho$};\draw[task axis,->] (0,-.2)--(0,3.9) node[above] {$p$};\node[below left] at (0,0) {$0$};\node[left] at (0,.6) {$p_0$};\draw[task curve] (.6,.6)--(2.4,.6);\draw[task curve,->] (1.3,.6)--(1.9,.6);\fill (.6,.6) circle (1.8pt) (2.4,.6) circle (1.8pt);\node[above right] at (.6,.6) {$1$};\node[above right] at (2.4,.6) {$2$};\node[below] at (.6,0) {$\rho_0$};\node[below] at (2.4,0) {$4\rho_0$};
\end{tikzpicture}\end{center}
Идеальный газ находится в сосуде под поршнем при температуре $800\,\text{К}$ и давлении $p_{0}=10^{5}\,\text{Па}$. На графике зависимости давления $p$ газа от его плотности $\rho$ изображен процесс перехода этого газа из состояния $1$ в состояние $2$. Определите температуру газа в состоянии $2$. (В кельвинах.)
}{
\smallskip\noindent\textbf{Решение}\par
Из уравнения состояния идеального газа: $pV = \frac{mRT}{M} .$ Учитывая, что $m/V = \rho,$ получаем: $p = \frac{\rho RT}{M} .$ 

В первом состоянии: $p_0 = \frac{\rho_0 RT_0}{M} ,$ во втором состоянии: $p_0 = \frac{4\rho_0 RT_2}{M} .$ Приравниваем правые части:

$\frac{\rho_0 RT_0}{M} = \frac{4\rho_0 RT_2}{M} .$

И получаем:

$T_2 = \frac{T_0}{4} = \frac{800}{4} = 200K.$

Ответ: $200$.
}

% Source sdamgia-43309: https://phys-ege.sdamgia.ru/problem?id=43309
\BankTask{sdamgia-43309}{
Температура идеального газа при постоянном давлении увеличилась в $2$ раза. При этом объем увеличился на $150$ $\text{дм}^{3}$. Чему равен начальный объем, занимаемый идеальным газом? (В кубических дециметрах.)
}{
\smallskip\noindent\textbf{Решение}\par
При изобарном процессе $\frac{V_1}{T_1} = \frac{V_2}{T_2} ,$ откуда следует, что при увеличении температуры в $2$ раза его объем увеличился в $2$ раза. Тогда получаем равенство $2V_1=V_1 + \Delta V,$ откуда первоначальный объем $V_1=\Delta V=150\text{дм} ^{3} .$

Ответ: $150$.
}

% Source sdamgia-31410: https://phys-ege.sdamgia.ru/problem?id=31410
\BankTask{sdamgia-31410}{
В закрытом сосуде неизменного объема при комнатной температуре находилась смесь аргона и гелия (по $1$ молю каждого газа). Половину содержимого этого сосуда выпустили наружу, а затем добавили в сосуд еще $2$ моля гелия. Газы считать идеальными, температуру — постоянной.

$1$. Парциальное давление аргона уменьшилось.

$2$. Давление смеси газов в сосуде не изменилось.

$3$. Концентрация гелия увеличилась.

$4$. В начале опыта плотности газов были одинаковые.

$5$. В начале опыта внутренние энергии газов были одинаковые.
}{
\smallskip\noindent\textbf{Решение}\par
$1$. Верно. Количество вещества аргона уменьшилось в $2$ раза, следовательно, из уравнения Клапейрона — Менделеева при неизменных объеме и температуре давление аргона уменьшилось.

$2$. Неверно. Вначале количество вещества в сосуде было равно $2$ моля, а после — $3$ моля. Следовательно, давление смеси газов увеличилось.

$3$. Верно. Поскольку количество вещества гелия увеличилось, то число молекул $N=\nu N_A$ также увеличилось, следовательно, концентрация $n= \frac{N}{V}$ увеличилась.

$4$. Неверно. Плотность вещества равна $\rho= \frac{m}{V} = \frac{\nu M}{V} .$ Поскольку молярные массы газов отличаются, то и плотность газов будет различной.

$5$. Верно. Внутренняя энергия одноатомных газов $U= \frac{3}{2} \nu RT$ будет одинаковой.

Ответ: $135$.
}

% Source sdamgia-44613: https://phys-ege.sdamgia.ru/problem?id=44613
\BankTask{sdamgia-44613}{
Объём сосуда, содержащего $1\,\text{моль}$ водорода, увеличили вдвое и добавили в сосуд $1{,}5$ моля гелия. Температуру в сосуде поддерживали постоянной.  

$1$) Концентрации водорода и гелия в сосуде в конце опыта одинаковы.

$2$) Внутренняя энергия водорода уменьшилась.

$3$) Плотность газа в сосуде увеличилась.

$4$) Давление в сосуде увеличилось.

$5$) Парциальное давление водорода уменьшилось.
}{
\smallskip\noindent\textbf{Решение}\par
$1$) Неверно. Число молекул гелия в $1{,}5$ раза больше, чем водорода, поэтому концентрация $n = \frac{N}{V}$ у газов будет различная.

$2$) Неверно. Внутренняя энергия водорода при неизменных температуре и количестве вещества не изменилась.

$3$) Верно. После добавления гелия в сосуд масса смеси газа стала равной 

$m = m_1 + m_2 = \nu_1 M_1 + \nu_2 M_2 = 1 \cdot 2 + 1{,}5 \cdot 4 = 2 + 6 = 8 \text{г},$

увеличилась в $4$ раза. Плотность при этом $\rho = \frac{m}{V}$ при увеличении массы в $4$ раза и объема в $2$ раза увеличилась в $2$ раза.

$4$) Верно. Из уравнения Клапейрона — Менделеева $p = \frac{\nu RT}{V}$ при увеличении количества вещества в $2{,}5$ раза и объема в $2$ раза при неизменной температуре давление смеси газов увеличилось.

$5$) Верно. Так как концентрация водорода при увеличении объема уменьшилась, то и давление водорода уменьшилось.

Ответ: $345$.
}

% Source sdamgia-6893: https://phys-ege.sdamgia.ru/problem?id=6893
\BankTask{sdamgia-6893}{
В цилиндре под поршнем находится идеальный одноатомный газ. Формулы $A$ и $\text{Б}$ ($p$ — давление; $V$ — объем; $\nu $ — количество вещества; $T$ — абсолютная температура) позволяют рассчитать значения физических величин, характеризующих состояние газа.  

\TaskMatchingColumns{Формулы}{А)  $\frac{\nu RT}{V}$\par  Б)  $\frac{\nu RT}{p}$}{Физические величины}{$1$)  давление\par $2$)  объем\par $3$)  молярная масса газа\par $4$)  масса газа}

\begin{center}\begin{tabular}{|c|c|}\hline
А & Б\\\hline
\strut & \strut\\\hline\end{tabular}\end{center}
}{
\smallskip\noindent\textbf{Решение}\par
Уравнение состояния идеального газа $pV_m=RT,$ где $V_m= \frac{V}{\nu}$ - молярный объем

Получаем формулу $p \frac{V}{\nu} =RT.$ Получаем $\frac{\nu RT}{V} =p$ и $\frac{\nu RT}{p} =V$

Ответ: $12$.
}

% Source sdamgia-1206: https://phys-ege.sdamgia.ru/problem?id=1206
\BankTask{sdamgia-1206}{
Воздух охлаждали в сосуде постоянного объема. При этом температура воздуха в сосуде снизилась в $4$ раза, а его давление уменьшилось в $2$ раза. Оказалось, что кран у сосуда был закрыт плохо и через него просачивался воздух. Во сколько раз увеличилась масса воздуха в сосуде?
}{
\smallskip\noindent\textbf{Решение}\par
В начальный момент времени воздух в стеклянном сосуде удовлетворял следующему уравнению состояния: $p_0V=\nu _0RT_0.$ В конечный момент — другому уравнению: $pV=\nu RT.$ При этом, согласно условию, $T= \frac{T_0}{4} ,$ а $p= \frac{p_0}{2} .$ Таким образом, $\frac{\nu}{\nu _0} =2,$ аналогичное соотношение верно и для массы газа. Следовательно, масса воздуха в сосуде увеличилась в $2$ раза.

Ответ: $2$.
}

% Source sdamgia-6433: https://phys-ege.sdamgia.ru/problem?id=6433
\BankTask{sdamgia-6433}{
Некоторое количество азота находится в замкнутом сосуде при давлении $1$ атм. Когда температуру сосуда повысили до $3000\,\text{К}$, давление увеличилось до $15$ атм, при этом половина имевшихся молекул азота распалась на атомы. Какой была температура газа до нагревания?
}{
\smallskip\noindent\textbf{Решение}\par
Воспользуемся формулой для давления газа: $p=nkT$. Если половина молекул азота распадается на атомы, то общее количество частиц в сосуде становится в полтора раза бóльшим. Во столько же раз возрастает их концентрация $n$. Согласно записанной формуле, если давление возросло в $15$ раз, а концентрация — в $1{,}5$ раза, то температура возросла в $10$ раз, и начальная температура была равна $300\,\text{К}$. 

Ответ: $300\,\text{К}$.
}

% Source sdamgia-11812: https://phys-ege.sdamgia.ru/problem?id=11812
\BankTask{sdamgia-11812}{
Два баллона объемами $10$ и $20\,\text{л}$ содержат $2\,\text{моль}$ кислорода и $1\,\text{моль}$ азота соответственно при температуре $28^\circ$C. Какое давление установится в баллонах, если их соединить между собой? Температуру газов считать неизменной.
}{
\smallskip\noindent\textbf{Решение}\par
В конечном состоянии в баллонах с суммарным объемом $V=30\text{л}$ будет находиться $\nu = 3\text{моль}$ смеси двух газов при температуре $28^\circ$C. Тогда согласно уравнению Менделеева — Клапейрона давление смеси будет:

$pV=\nu RT .$

$p= \frac{\nu RT}{V} = \frac{3 \cdot 8{,}31 \cdot 301}{30 \cdot 10 ^{- 3}} \approx 250 \cdot 10 ^{3} \text{Па}=250\text{кПа}.$

Ответ: $250\,\text{кПа}$.
}


% Source 772AAE: https://ege.fipi.ru/bank/questions.php?search=1&proj=BA1F39653304A5B041B656915DC36B38&qid=772AAE
\BankTask{772AAE}{
В сосуде содержится разреженный аргон, абсолютная температура которого равна \(150\ \text{К}\). Концентрацию аргона уменьшили в \(2\) раза, при этом его давление увеличилось в \(3\) раза. Определите абсолютную температуру газа в конечном равновесном состоянии. (в $\text{К}$).
}{
\smallskip\noindent\textbf{Решение от GPT}\par
Из уравнения \(p=nkT\) следует
\[
\frac{T_2}{T_1}=\frac{p_2}{p_1}\frac{n_1}{n_2}=3\cdot2=6.
\]
Следовательно, \(T_2=6\cdot150=900\ \text{К}\).
Ответ: \(900\ \text{К}\).
}

% Source AD774B: https://ege.fipi.ru/bank/questions.php?search=1&proj=BA1F39653304A5B041B656915DC36B38&qid=AD774B
\BankTask{AD774B}{
В сосуде неизменного объёма находится идеальный газ. Во сколько раз нужно увеличить количество газа в сосуде, чтобы после уменьшения абсолютной температуры газа в \(2\) раза его давление стало вдвое больше начального?
}{
\smallskip\noindent\textbf{Решение от GPT}\par
По уравнению \(pV=\nu RT\) при неизменном объёме
\[
\frac{\nu_2}{\nu_1}=\frac{p_2}{p_1}\frac{T_1}{T_2}=2\cdot2=4.
\]
Ответ: в \(4\) раза.
}

% Source 8E16FE: https://ege.fipi.ru/bank/questions.php?search=1&proj=BA1F39653304A5B041B656915DC36B38&qid=8E16FE
\BankTask{8E16FE}{
Объём \(1\ \text{моль}\) водорода в сосуде при температуре \(T_0\) и давлении \(p_0\) равен \(20\ \text{л}\). Каков объём \(3\ \text{моль}\) водорода при том же давлении и вдвое большей температуре (в $\text{л}$)?
}{
\smallskip\noindent\textbf{Решение от GPT}\par
Из уравнения \(pV=\nu RT\) при одинаковом давлении получаем
\[
\frac{V_2}{V_1}=\frac{\nu_2}{\nu_1}\frac{T_2}{T_1}=3\cdot2=6.
\]
Тогда \(V_2=6\cdot20=120\ \text{л}\).
Ответ: \(120\ \text{л}\).
}

% Source 5F4252: https://ege.fipi.ru/bank/questions.php?search=1&proj=BA1F39653304A5B041B656915DC36B38&qid=5F4252
\BankTask{5F4252}{
\begin{center}
\begin{tikzpicture}[x=1cm,y=0.065cm,>=stealth]
\draw[->] (0,0)--(4.7,0) node[right] {$V,\,\text{м}^{3}$};
\draw[->] (0,0)--(0,59) node[left] {$p,\,\text{кПа}$};
\node[below left] at (0,0) {$0$};
\foreach \x in {1,2,3} {\draw (\x,-2)--(\x,2);}
\node[below] at (1,-2) {$1$};
\node[below] at (3,-2) {$3$};
\foreach \y in {20,40} {\draw (-0.12,\y)--(1.25,\y);\node[left] at (-0.12,\y) {$\y$};}
\draw (2.5,20)--(3,20);
\draw (2.5,40)--(3,40);
\draw[thick] (3,20)--(3,40);
\fill (3,20) circle (1.5pt) node[right] {$1$};
\fill (3,40) circle (1.5pt) node[right] {$2$};
\end{tikzpicture}
\end{center}
На $pV$-диаграмме представлен процесс $1$--$2$, в котором участвует разреженный аргон ($p$ — давление аргона, $V$ — его объём). Абсолютная температура газа в состоянии $1$ равна $250\,\text{К}$. Какая абсолютная температура соответствует состоянию $2$, если масса газа остаётся неизменной (в $\text{К}$)?
}{
\smallskip\noindent\textbf{Решение от GPT}\par
При неизменной массе газа $pV=\nu RT$. Объём постоянен, поэтому
\[
T_2=T_1\frac{p_2}{p_1}=250\cdot\frac{40}{20}=500\,\text{К}.
\]
Ответ: $500\,\text{К}$.
}

% Source B86966: https://ege.fipi.ru/bank/questions.php?search=1&proj=BA1F39653304A5B041B656915DC36B38&qid=B86966
\BankTask{B86966}{
\begin{center}
\begin{tikzpicture}[x=0.012cm,y=0.8cm,>=stealth]
\draw[->] (0,0)--(530,0) node[right] {$T,\,\text{К}$};
\draw[->] (0,0)--(0,5) node[left] {$V,\,\text{м}^{3}$};
\node[below left] at (0,0) {$0$};
\draw (273,-0.08)--(273,0.08);
\node[below] at (273,-0.08) {$273$};
\draw (35,1)--(210,1);
\draw (65,4)--(245,4);
\node[left] at (0,1) {$1$};
\node[left] at (0,4) {$4$};
\draw[dashed] (273,0)--(273,1);
\draw (273,1)--(273,4);
\draw[->,thick] (273,1)--(273,2.5);
\fill (273,1) circle (1.3pt) node[right] {$1$};
\fill (273,4) circle (1.3pt) node[right] {$2$};
\end{tikzpicture}
\end{center}
На рисунке представлен график зависимости объёма идеального газа от его температуры в некотором процессе. В состоянии $1$ давление газа было равно нормальному атмосферному давлению. Какое давление соответствует состоянию $2$, если масса газа остаётся неизменной (в $\text{кПа}$)?
}{
\smallskip\noindent\textbf{Решение от GPT}\par
По графику $T_1=T_2=273\,\text{К}$, $V_1=1\,\text{м}^{3}$, $V_2=4\,\text{м}^{3}$. Поэтому $p_1V_1=p_2V_2$. Принимая нормальное атмосферное давление $p_1\approx100\,\text{кПа}$, получаем
\[
p_2=p_1\frac{V_1}{V_2}=100\cdot\frac{1}{4}=25\,\text{кПа}.
\]
Ответ: $25\,\text{кПа}$.
}

% Source 79B504: https://ege.fipi.ru/bank/questions.php?search=1&proj=BA1F39653304A5B041B656915DC36B38&qid=79B504
\BankTask{79B504}{
\begin{center}
\begin{tikzpicture}[x=0.11cm,y=5.5cm,>=stealth]
\draw[thin,gray] (0,0) grid[xstep=10,ystep=0.2] (40,0.6);
\draw[->] (0,0)--(51,0) node[right] {$T,\,\text{К}$};
\draw[->] (0,0)--(0,0.82) node[left] {$p,\,10^{4}\,\text{Па}$};
\foreach \x/\lab in {0/290,10/300,20/310,30/320,40/330} {
\draw (\x,0)--(\x,-0.015);
\node[below] at (\x,-0.015) {$\lab$};
}
\foreach \y/\lab in {0/{1{,}8},0.2/{2{,}0},0.4/{2{,}2},0.6/{2{,}4}} {
\draw (-0.7,\y)--(0,\y);
\node[left] at (-0.7,\y) {$\lab$};
}
\draw[thick] (0,0.1333)--(44,0.4267);
\end{tikzpicture}
\end{center}
На рисунке показан график зависимости давления газа в запаянном сосуде от его температуры. Объём сосуда равен $0{,}25\,\text{м}^{3}$. Какое приблизительно количество газообразного вещества содержится в этом сосуде (в $\text{моль}$, с округлением до целого)?
}{
\smallskip\noindent\textbf{Решение от GPT}\par
По графику при $T=300\,\text{К}$ давление $p\approx2{,}0\cdot10^{4}\,\text{Па}$. Из уравнения идеального газа
\[
\nu=\frac{pV}{RT}=\frac{2{,}0\cdot10^{4}\cdot0{,}25}{8{,}31\cdot300}\approx2{,}01\,\text{моль}.
\]
После округления ответ: $2\,\text{моль}$.
}

% Source A1DD9C: https://ege.fipi.ru/bank/questions.php?search=1&proj=BA1F39653304A5B041B656915DC36B38&qid=A1DD9C
\BankTask{A1DD9C}{
Из металлического баллона, в котором находилось $3\,\text{моль}$ идеального газа под давлением $3\cdot10^5\,\text{Па}$, медленно, без изменения температуры выпустили $0{,}6\,\text{моль}$ газа. Определите давление газа в конечном состоянии. (в $\text{кПа}$).
}{
\smallskip\noindent\textbf{Решение от GPT}\par
При постоянных $T$ и $V$ давление пропорционально количеству вещества: $p_2=p_1\nu_2/\nu_1=3\cdot10^5(3-0{,}6)/3=2{,}4\cdot10^5\,\text{Па}$. Ответ: $240\,\text{кПа}$.
}

% Source E27EA7: https://ege.fipi.ru/bank/questions.php?search=1&proj=BA1F39653304A5B041B656915DC36B38&qid=E27EA7
\BankTask{E27EA7}{
При изучении процессов, происходящих с гелием, ученик занёс в таблицу результаты измерения температуры и давления одного и того же количества газа в различных равновесных состояниях. Газ считать идеальным.
\begin{center}
\begin{tabular}{|c|c|c|c|c|c|c|c|}
\hline
№ состояния & 1 & 2 & 3 & 4 & 5 & 6 & 7 \\
\hline
$p,\,\text{кПа}$ & $100$ & $90$ & $75$ & $50$ & $55$ & $75$ & $100$ \\
\hline
$t$, ${}^{\circ}\mathrm{C}$ & $27$ & $27$ & $27$ & $27$ & $57$ & $177$ & $327$ \\
\hline
\end{tabular}
\end{center}
\begin{enumerate}
\item Объём газа в состоянии $4$ в $2$ раза меньше объёма газа в состоянии $1$.
\item В состояниях $4$--$7$ объём газа был одинаковым.
\item Внутренняя энергия газа в состоянии $6$ в $3$ раза больше, чем в состоянии $5$.
\item При переходе от состояния $2$ к состоянию $3$ в ходе изотермического процесса газ получал тепло.
\item При переходе от состояния $5$ к состоянию $6$ в ходе изохорного процесса газ совершал работу.
\end{enumerate}
}{
\smallskip\noindent\textbf{Решение от GPT}\par
Для данного количества газа $V=\nu RT/p$, а для гелия $U=\frac32\nu RT$. Температуры состояний: $300$, $300$, $300$, $300$, $330$, $450$, $600\ \text{К}$.
\begin{enumerate}
\item Неверно: $V_4/V_1=p_1/p_4=2$, то есть объём больше в $2$ раза.
\item Верно: $T_4/p_4=T_5/p_5=T_6/p_6=T_7/p_7=6\ \text{К}/\text{кПа}$, поэтому объёмы одинаковы.
\item Неверно: $U_6/U_5=T_6/T_5=450/330=15/11\ne3$.
\item Верно: при постоянной температуре давление уменьшается, газ расширяется; $\Delta U=0$, $Q=A>0$.
\item Неверно: при изохорном процессе $A=0$.
\end{enumerate}
Ответ: $24$.
}

% Source 84BDA5: https://ege.fipi.ru/bank/questions.php?search=1&proj=BA1F39653304A5B041B656915DC36B38&qid=84BDA5
\BankTask{84BDA5}{
Сосуд разделён на две равные по объёму части пористой неподвижной перегородкой. В левой части сосуда содержится $8\ \text{г}$ гелия, в правой --- $1\ \text{моль}$ аргона. Перегородка может пропускать молекулы гелия и является непроницаемой для молекул аргона. Температура газов одинакова и остаётся постоянной.
Рассматривается состояние газов после установления равновесия в системе.
\begin{enumerate}
\item Давление в обеих частях сосуда одинаково.
\item Концентрация гелия и аргона в правой части сосуда одинакова.
\item Внутренняя энергия гелия в сосуде больше, чем внутренняя энергия аргона.
\item Внутренняя энергия гелия в сосуде в конечном состоянии больше, чем в начальном.
\item В правой части сосуда общее число молекул газов в $2$ раза меньше, чем в левой части.
\end{enumerate}
}{
\smallskip\noindent\textbf{Решение от GPT}\par
Молярная масса гелия $M_{\mathrm{He}}=4\ \text{г/моль}$, поэтому $\nu_{\mathrm{He}}=8/4=2\ \text{моль}$. В равновесии концентрация гелия одинакова в обеих частях: в каждой находится по $1\ \text{моль}$ гелия. Весь аргон ($1\ \text{моль}$) остаётся справа. Пусть объём каждой части равен $V$.
\begin{enumerate}
\item Неверно: $p_{\text{лев}}=\dfrac{(1\ \text{моль})RT}{V}$, $p_{\text{прав}}=2p_{\text{лев}}$.
\item Верно: справа количества гелия и аргона равны, поэтому равны и их концентрации.
\item Верно: оба газа одноатомные, $U=\frac32\nu RT$; следовательно, $U_{\mathrm{He}}/U_{\mathrm{Ar}}=2$.
\item Неверно: количество гелия и его температура не изменились, поэтому его внутренняя энергия постоянна.
\item Неверно: справа общее число молекул в $2$ раза больше, чем слева.
\end{enumerate}
Ответ: $23$.
}

% Source 403D76: https://ege.fipi.ru/bank/questions.php?search=1&proj=BA1F39653304A5B041B656915DC36B38&qid=403D76
\BankTask{403D76}{
В двух сосудах одинакового объёма находятся разреженные газы.
В первом сосуде находится \(2\,\text{моль}\) гелия при температуре \(127\,^{\circ}\mathrm{C}\), во втором сосуде находится \(1\,\text{моль}\) аргона при температуре \(300\,\text{К}\).
\begin{enumerate}
\item Абсолютная температура газа во втором сосуде выше, чем в первом.
\item Среднеквадратичная скорость молекул газа в первом сосуде больше, чем во втором.
\item Давления газов в сосудах одинаковы.
\item Концентрация газа в первом сосуде в \(2\) раза меньше, чем во втором.
\item Отношение средней кинетической энергии теплового движения молекул аргона к средней кинетической энергии теплового движения молекул гелия равно \(0{,}75\).
\end{enumerate}
}{
\smallskip\noindent\textbf{Решение от GPT}\par
\(T_1=127+273=400\,\text{К}\), \(T_2=300\,\text{К}\).
\begin{enumerate}
\item Неверно: \(T_2<T_1\).
\item Верно: \(v_{\text{кв}}=\sqrt{3RT/M}\); у гелия температура выше, а молярная масса меньше, поэтому \(v_{\text{кв},1}>v_{\text{кв},2}\).
\item Неверно: \(p_1/p_2=\nu_1T_1/(\nu_2T_2)=8/3\).
\item Неверно: при одинаковых объёмах \(n_1/n_2=\nu_1/\nu_2=2\), то есть концентрация в первом сосуде больше.
\item Верно: \(\overline E=3kT/2\), поэтому \(\overline E_{\mathrm{Ar}}/\overline E_{\mathrm{He}}=300/400=0{,}75\).
\end{enumerate}
Ответ: \(25\).
}

% Source 2CEC9C: https://ege.fipi.ru/bank/questions.php?search=1&proj=BA1F39653304A5B041B656915DC36B38&qid=2CEC9C
\BankTask{2CEC9C}{
Объём сосуда, содержащего \(1\,\text{моль}\) неона, уменьшили вдвое и добавили в сосуд \(1\,\text{моль}\) гелия. Температура в сосуде поддерживается постоянной.
\begin{enumerate}
\item Концентрации неона и гелия в сосуде одинаковы.
\item Внутренняя энергия неона уменьшилась.
\item Плотность газа в сосуде не изменилась.
\item Давление в сосуде увеличилось.
\item Парциальное давление неона не изменилось.
\end{enumerate}
}{
\smallskip\noindent\textbf{Решение от GPT}\par
Пусть начальные объём и давление равны \(V_0\) и \(p_0\).

1) Верно: количества неона и гелия одинаковы, они занимают общий объём, поэтому \(n_{\mathrm{Ne}}=n_{\mathrm{He}}\).

2) Неверно: \(U_{\mathrm{Ne}}=\frac32\nu_{\mathrm{Ne}}RT\) не изменяется.

3) Неверно: масса газа увеличилась, а объём уменьшился вдвое, поэтому плотность увеличилась.

4) Верно: по уравнению идеального газа \(p\propto\nu/V\), поэтому \(p_f/p_0=2\cdot2=4\).

5) Неверно: \(p_{\mathrm{Ne}}=\nu_{\mathrm{Ne}}RT/V\), поэтому парциальное давление неона увеличилось вдвое.

Ответ: \(14\).
}

% Source CD1D80: https://ege.fipi.ru/bank/questions.php?search=1&proj=BA1F39653304A5B041B656915DC36B38&qid=CD1D80
\BankTask{CD1D80}{
В цилиндре под поршнем находится идеальный одноатомный газ. Формулы А и Б позволяют рассчитать значения физических величин, характеризующих состояние газа. Обозначения: \(p\) — давление; \(T\) — абсолютная температура; \(N\) — число атомов газа; \(k\) — постоянная Больцмана, \(V\) — объём газа.

Установите соответствие между формулами и физическими величинами, значение которых можно рассчитать по этим формулам.

\begin{center}
\TaskMatchingColumns[.35\linewidth]{Формулы}{
А) \(\displaystyle\frac{p}{kT}\)\par
Б) \(\displaystyle\frac{pV}{Nk}\)
}{Физические величины}{
1) абсолютная температура\par
2) концентрация молекул\par
3) давление\par
4) внутренняя энергия
}

\begin{tabular}{|c|c|}
\hline
А & Б \\
\hline
\phantom{0} & \phantom{0} \\
\hline
\end{tabular}
\end{center}
}{
\smallskip\noindent\textbf{Решение от GPT}\par
Из \(p=nkT\) следует \(n=\frac{p}{kT}\), то есть формула А определяет концентрацию молекул: А — 2. Из \(pV=NkT\) получаем \(T=\frac{pV}{Nk}\), то есть формула Б определяет абсолютную температуру: Б — 1.

Ответ: А — 2, Б — 1; \(21\).
}

% Source F5A444: https://ege.fipi.ru/bank/questions.php?search=1&proj=BA1F39653304A5B041B656915DC36B38&qid=F5A444
\BankTask{F5A444}{
Один моль одноатомного идеального газа участвует в процессе $1$--$2$--$3$, график которого изображён на рисунке в координатах $p$--$T$, где $p$~--- давление газа, $T$~--- абсолютная температура газа. Как изменяются объём газа $V$ в ходе процесса $1$--$2$ и плотность газа $\rho$ в ходе процесса $2$--$3$?

\begin{center}
\begin{tikzpicture}[x=1cm,y=1cm,>=stealth]
\draw[->] (0,0)--(4,0) node[below] {$T$};
\draw[->] (0,0)--(0,3.3) node[left] {$p$};
\node[below left] at (0,0) {$0$};
\draw (0,0)--(2.9,1.22)--(2.9,3.05);
\fill (1.03,0.433) circle (1.2pt) node[above left] {$1$};
\fill (2.9,1.22) circle (1.2pt) node[right] {$2$};
\fill (2.9,3.05) circle (1.2pt) node[right] {$3$};
\draw[->] (1.4,0.589)--(1.95,0.82);
\draw[->] (2.9,1.7)--(2.9,2.2);
\end{tikzpicture}
\end{center}
Для каждой величины определите соответствующий характер изменения:

\begin{enumerate}
\item увеличивается
\item уменьшается
\item не изменяется
\end{enumerate}


\TaskChangeTable{Объём газа в ходе процесса $1$--$2$}{Плотность газа в ходе процесса $2$--$3$}
}{
\smallskip\noindent\textbf{Решение от GPT}\par
Из $pV=\nu RT$ следует $V=\nu RT/p$. На участке $1$--$2$ прямая проходит через начало координат, поэтому $p/T=\mathrm{const}$ и объём не изменяется (3).

На участке $2$--$3$ температура постоянна, давление увеличивается, поэтому объём уменьшается. При постоянной массе $\rho=m/V$ увеличивается (1).

Ответ: $31$.
}

% Source F0D87A: https://ege.fipi.ru/bank/questions.php?search=1&proj=BA1F39653304A5B041B656915DC36B38&qid=F0D87A
\BankTask{F0D87A}{
$1$ моль одноатомного идеального газа участвует в процессе $1$--$2$--$3$, график которого изображён на рисунке в координатах $p$--$V$, где $p$~--- давление газа, $V$~--- объём газа. Как изменяются концентрация молекул $n$ газа в ходе процесса $1$--$2$ и абсолютная температура $T$ газа в ходе процесса $2$--$3$? Масса газа остаётся постоянной.

\begin{center}
\begin{tikzpicture}[x=1cm,y=1cm,>=stealth]
\draw[->] (0,0)--(4.2,0) node[below] {$V$};
\draw[->] (0,0)--(0,3.5) node[left] {$p$};
\node[below left] at (0,0) {$0$};
\draw (0.8,0.9)--(3.2,0.9)--(3.2,2.9);
\fill (0.8,0.9) circle (1.2pt) node[below left] {$1$};
\fill (3.2,0.9) circle (1.2pt) node[below right] {$2$};
\fill (3.2,2.9) circle (1.2pt) node[above right] {$3$};
\draw[->] (1.4,0.9)--(2,0.9);
\draw[->] (3.2,1.4)--(3.2,1.95);
\end{tikzpicture}
\end{center}
Для каждой величины определите соответствующий характер изменения:

\begin{enumerate}
\item увеличивается
\item уменьшается
\item не изменяется
\end{enumerate}

\TaskChangeTable{Концентрация молекул газа в ходе процесса $1$--$2$}{Абсолютная температура газа в ходе процесса $2$--$3$}
}{
\smallskip\noindent\textbf{Решение от GPT}\par
На участке $1$--$2$ объём увеличивается. Число молекул $N$ постоянно, поэтому концентрация $n=N/V$ уменьшается (2).

На участке $2$--$3$ объём постоянен, давление увеличивается. Из $T=pV/(\nu R)$ следует, что температура увеличивается (1).

Ответ: $21$.
}

% Source F12DcD: https://ege.fipi.ru/bank/questions.php?search=1&proj=BA1F39653304A5B041B656915DC36B38&qid=F12DcD
\BankTask{F12DcD}{
\begin{center}
\begin{tikzpicture}[>=stealth,scale=0.85]
\draw[->] (0,0)--(4,0) node[below] {$T$};
\draw[->] (0,0)--(0,3.8) node[left] {$V$};
\node[below left] at (0,0) {$0$};
\coordinate (A) at (2.5,3.2);
\coordinate (B) at (2.5,1);
\coordinate (C) at (1.3,0.52);
\draw[dashed] (0,0)--(C);
\draw (A)--(B)--(C);
\draw[->] (2.5,2.65)--(2.5,1.95);
\draw[->] (2.25,0.9)--(1.75,0.7);
\fill (A) circle (1.5pt) node[right] {$1$};
\fill (B) circle (1.5pt) node[right] {$2$};
\fill (C) circle (1.5pt) node[above left] {$3$};
\end{tikzpicture}
\end{center}
Один моль идеального газа участвует в процессе \(1\text{–}2\text{–}3\), график которого изображён на рисунке в координатах \(V\text{–}T\), где \(V\) — объём газа, \(T\) — абсолютная температура газа. Как изменяются давление газа \(p\) в ходе процесса \(1\text{–}2\) и концентрация молекул газа \(n\) в ходе процесса \(2\text{–}3\)? Масса газа остаётся постоянной.

Для каждой величины определите соответствующий характер изменения:
\begin{enumerate}
\item увеличивается
\item уменьшается
\item не изменяется
\end{enumerate}
\TaskChangeTable{Давление газа в ходе процесса \(1\text{–}2\)}{Концентрация молекул газа в ходе процесса \(2\text{–}3\)}
}{
\smallskip\noindent\textbf{Решение от GPT}\par
В процессе \(1\text{–}2\) температура постоянна, а объём уменьшается. Из \(p=\nu RT/V\) следует, что давление увеличивается (\(1\)). В процессе \(2\text{–}3\) объём уменьшается при постоянном числе молекул, поэтому концентрация \(n=N/V\) увеличивается (\(1\)).

Ответ: \(11\).
}

% Source C427F7: https://ege.fipi.ru/bank/questions.php?search=1&proj=BA1F39653304A5B041B656915DC36B38&qid=C427F7
\BankTask{C427F7}{
Объём сосуда с идеальным газом уменьшили вдвое, выпустив половину газа и поддерживая температуру газа в сосуде постоянной. Как изменились при этом давление газа в сосуде и его внутренняя энергия?

Для каждой величины определите соответствующий характер изменения:

\begin{enumerate}
\item увеличилась
\item уменьшилась
\item не изменилась
\end{enumerate}


\TaskChangeTable{Давление газа в сосуде}{Внутренняя энергия газа в сосуде}
}{
\smallskip\noindent\textbf{Решение от GPT}\par
При постоянной температуре \(p=\nu RT/V\). Поскольку \(\nu_2=\nu_1/2\) и \(V_2=V_1/2\), получаем \(p_2=p_1\): давление не изменилось (\(3\)). Внутренняя энергия идеального газа при неизменной температуре пропорциональна количеству вещества: \(U_2=U_1/2\), поэтому она уменьшилась (\(2\)).

Ответ: \(32\).
}

% Source FB5B0D: https://ege.fipi.ru/bank/questions.php?search=1&proj=BA1F39653304A5B041B656915DC36B38&qid=FB5B0D
\BankTask{FB5B0D}{
В сосуде неизменного объёма при комнатной температуре находилась смесь неона и аргона, по \(2\,\text{моль}\) каждого. Половину смеси выпустили, а затем добавили в сосуд \(2\,\text{моль}\) аргона. Как в результате этого изменились парциальное давление аргона и давление смеси газов, если температура смеси оставалась неизменной? Для каждой величины определите соответствующий характер изменения:

\begin{enumerate}
\item увеличилась
\item уменьшилась
\item не изменилась
\end{enumerate}


\TaskChangeTable{Парциальное давление аргона}{Давление смеси газов в сосуде}
}{
\smallskip\noindent\textbf{Решение от GPT}\par
После выпуска половины смеси осталось по \(1\,\text{моль}\) каждого газа. После добавления аргона \(\nu_{\mathrm{Ar},2}=3\,\text{моль}\), а общее количество вещества \(\nu_2=4\,\text{моль}\). При постоянных \(T\) и \(V\) давление пропорционально количеству вещества. Поэтому парциальное давление аргона увеличилось в \(3/2\) раза (\(1\)), а давление смеси не изменилось, поскольку первоначально также было \(4\,\text{моль}\) газа (\(3\)).

Ответ: \(13\).
}

% Source 3BDF8F: https://ege.fipi.ru/bank/questions.php?search=1&proj=BA1F39653304A5B041B656915DC36B38&qid=3BDF8F
\BankTask{3BDF8F}{
В сосуде постоянного объёма абсолютную температуру одноатомного идеального газа увеличили в \(3\) раза, выпустив при этом \(\frac{1}{3}\) газа из сосуда. Как изменились в результате этого плотность газа в сосуде и его внутренняя энергия?

Для каждой величины определите соответствующий характер изменения:
\begin{enumerate}
\item увеличилась
\item уменьшилась
\item не изменилась
\end{enumerate}

\TaskChangeTable{Плотность газа в сосуде}{Внутренняя энергия газа в сосуде}
}{
\smallskip\noindent\textbf{Решение от GPT}\par
При постоянном объёме \(\rho=m/V\). После выпуска газа \(m_2=\frac{2}{3}m_1\), поэтому \(\rho_2=\frac{2}{3}\rho_1\): плотность уменьшилась (\(2\)).

Для одноатомного идеального газа \(U=\frac{3}{2}\nu RT\), поэтому
\[
\frac{U_2}{U_1}=\frac{\nu_2}{\nu_1}\frac{T_2}{T_1}=\frac{2}{3}\cdot3=2.
\]
Внутренняя энергия увеличилась (\(1\)). Ответ: \(21\).
}

% Source cBe7DB: https://ege.fipi.ru/bank/questions.php?search=1&proj=BA1F39653304A5B041B656915DC36B38&qid=cBe7DB
\BankTask{cBe7DB}{
Ученику необходимо на опыте обнаружить зависимость объёма газа, находящегося в сосуде под подвижным поршнем, от молярной массы газа.

У него имеется пять различных сосудов с манометрами и термометрами. Сосуды наполнены равными массами различных газов при различных температурах (см. таблицу). Какие два сосуда необходимо взять ученику, чтобы провести исследование?

\begin{center}
\begin{tabular}{|c|c|c|c|}
\hline
\shortstack{№\\сосуда} & Давление, \(\text{кПа}\) & \shortstack{Температура газа\\в сосуде, \(\text{К}\)} & Газ в сосуде \\
\hline
\(1\) & \(100\) & \(280\) & кислород \\
\hline
\(2\) & \(50\) & \(270\) & кислород \\
\hline
\(3\) & \(100\) & \(280\) & азот \\
\hline
\(4\) & \(50\) & \(300\) & азот \\
\hline
\(5\) & \(60\) & \(320\) & кислород \\
\hline
\end{tabular}
\end{center}
}{
\smallskip\noindent\textbf{Решение от GPT}\par
Из уравнения состояния идеального газа \(V=\dfrac{mRT}{Mp}\): для исследования зависимости объёма от молярной массы \(M\) нужно сохранить одинаковыми массу, температуру и давление газа. Массы всех газов равны, а в сосудах \(1\) и \(3\) одинаковы давление \(100\,\text{кПа}\) и температура \(280\,\text{К}\), но различны газы и их молярные массы. Ответ: \(13\).
}

% Source 49644: https://phys-ege.sdamgia.ru/problem?id=49644
\BankTask{49644}{
Метеорологический зонд общей массой $20\,\text{кг}$ удерживают на поверхности Земли силой, равной по модулю $1$ кН и направленной вниз. Затем зонд отпускают, он поднимается вверх и остаётся на такой высоте, где его объём увеличивается в $2$ раза. Температура, измеренная зондом на этой высоте, равна –$43$ \ensuremath{{}^\circ}С. Определите давление на этой высоте, если на поверхности Земли давление равно $10^{5}$ Па, а температура равна +$17$ \ensuremath{{}^\circ}С.
}{
\smallskip\noindent\textbf{Решение}\par
На зонд, находящийся на Земле, действуют сила тяжести mg,  приложенная сила F, направленные вертикально вниз, сила Архимеда $F_{\,\text{арх}1,}$ направленная вертикально вверх. Зонд неподвижный, поэтому по второму закону Ньютона

\[F_{A1} = mg + F.\]

Сила Архимеда в этом случае равна $F_{A1}=\rho_1gV,$ где $\rho_1$ — плотность воздуха у поверхности Земли, V  — первоначальный объем зонда.
Плотность воздуха  найдем из уравнения Менделеева  — Клайперона $p_1 = \frac{\rho_1RT_1}{M} ,$ откуда $\rho_1 = \frac{p_1M}{RT_1} .$
Когда зонд поднимется на некоторую высоту, на него будут действовать сила тяжести и сила Архимеда $F_{\,\text{арх}}2 = \rho_2g \cdot 2V.$ Плотность воздуха на высоте находим аналогично из уравнения Менделеева  — Клапейрона $\rho_2 = \frac{p_2M}{RT_2} .$
Решая систему уравнений, находим давление на высоте:

\[p_2 = \frac{mgp_1T_2}{2T_1 ( mg + F )} = \frac{20 \cdot 10 \cdot 10 ^{5} \cdot 230}{2 \cdot 290 \cdot ( 20 \cdot 10 + 1000 )} \approx 6{,}6\,\text{кПа}.\]

Ответ: $6{,}6$ кПа.
}

% Source 1210: https://phys-ege.sdamgia.ru/problem?id=1210
\BankTask{1210}{
Из стеклянного сосуда стали выпускать сжатый воздух, одновременно охлаждая сосуд. При этом температура воздуха упала вдвое, а его давление уменьшилось в $3$ раза. Во сколько раз уменьшилась масса воздуха?
}{
\smallskip\noindent\textbf{Решение}\par
В начальный момент времени воздух в стеклянном сосуде удовлетворял следующему уравнению состояния: $p_0V=\nu _0RT_0.$ В конечный момент  — другому уравнению: $pV=\nu RT.$ При этом, согласно условию, $T= \frac{T_0}{2} ,$ а $p= \frac{p_0}{3} .$ Таким образом:

\[\frac{\nu}{\nu _0} = \frac{2}{3} ,\]
 аналогичное соотношение верно и для массы газа. Следовательно, масса воздуха в сосуде уменьшилась в $1{,}5$ раза.

Ответ: $1{,}5$.
}

% Source 2980: https://phys-ege.sdamgia.ru/problem?id=2980
\BankTask{2980}{
В сосуде с небольшой трещиной находится воздух. Воздух может медленно просачиваться сквозь трещину. Во время опыта объем сосуда уменьшили в $8$ раз, давление воздуха в сосуде увеличилось в $2$ раза, а его абсолютная температура увеличилась в $1{,}5$ раза. Во сколько раз изменилась внутренняя энергия воздуха в сосуде? (Воздух считать идеальным газом.)
}{
\smallskip\noindent\textbf{Решение}\par
Внутренняя энергия газа пропорциональна его температуре и количеству вещества (газа) в сосуде: $U= \frac{i}{2} \nu R T,$ где для воздуха можно принять $i=5.$ Согласно уравнению Клапейрона  — Менделеева, $pV=\nu RT.$ Следовательно, $U= \frac{i}{2} pV.$ Согласно условию задачи, p возросло в $2$ раза, а V уменьшилось в $8$ раз. Следовательно, U уменьшилась в $4$ раза.

Ответ: внутренняя энергия воздуха в сосуде уменьшилась в $4$ раза.
}

% Source 19678: https://phys-ege.sdamgia.ru/problem?id=19678
\BankTask{19678}{
В закрытом сосуде с клапаном находится идеальный газ при температуре +$20$ °С и давлении p$_{1}$. В результате некоторого эксперимента $20$ \% газа вышло из сосуда через клапан. При этом температура газа повысилась на $10$ °С, а его давление изменилось до некоторой величины p$_{2}$. Найдите отношение $\frac{p_1}{p_2} .$ Ответ округлите до десятых долей.
}{
\smallskip\noindent\textbf{Решение}\par
Количество газа в сосуде, оставшегося после некоторого эксперимента, составляет $0{,}8$ от изначального количества $\nu_2=0{,}8\nu_1.$ Температура после эксперимента составит $303$ К. Из уравнения Менделеева—Клапейрона $pV=\nu RT$ получим: 

\[\frac{p_1}{p_2} = \frac{\nu_1 RT_1}{\nu_2 RT_2} = \frac{\nu_1T_1}{0{,}8\nu_1T_2} = \frac{293}{0{,}8 \cdot 303} \approx 1{,}2.\]

Ответ: $1{,}2$.
}

% Source 32562: https://phys-ege.sdamgia.ru/problem?id=32562
\BankTask{32562}{
В закрытом сосуде объемом V  =  $30$ л находится влажный воздух массой m  =  $45\,\text{г}$ при температуре t  =  $80$ °C и давлении p  =  $160$ кПа. Определите относительную влажность воздуха, находящегося в сосуде. Давление $p_{\,\text{н}}$ насыщенных паров воды при температуре $80$ °С равно $47$ кПа.
}{
\smallskip\noindent\textbf{Решение}\par
Относительная влажность воздуха определяется по формуле $\varphi = \frac{p_{\,\text{п}}}{p_{\,\text{н}}} \cdot 100\%.$
По закону Дальтона давление в сосуде складывается из парциального давления водяного пара $p_{\,\text{п}}$ и давления сухого воздуха $p_{\,\text{возд}}$, то есть $p=p_{\,\text{п}} + p_{\,\text{возд}}.$
Из уравнения Клапейрона  — Менделеева найдем давление сухого воздуха:

$p_{\,\text{возд}}V= \frac{m_{\,\text{возд}}RT}{M_{\,\text{возд}}} ,$ следовательно,  $p_{\,\text{возд}}= \frac{m_{\,\text{возд}}RT}{VM_{\,\text{воз}}} .$

Аналогично находим парциальное давление водяного пара:

$p_{\,\text{п}}V= \frac{m_{\,\text{п}}RT}{M_{\,\text{вод}}} ,$ следовательно, $p_{\,\text{п}}= \frac{m_{\,\text{п}}RT}{VM_{\,\text{вод}}} .$

Учтем, что масса влажного воздуха равна $m=m_{\,\text{п}} + m_{\,\text{возд}}.$ Объединяя уравнения, выразим парциальное давление водяного пара:

\[p_{\,\text{п}}= \frac{pVM_{\,\text{возд}} - mRT}{V ( M_{\,\text{возд}} - M_{\,\text{вод}} )} .\]

Тогда относительная влажность воздуха равна:

\[\varphi = \frac{pVM_{\,\text{возд}} - mRT}{p_{\,\text{н}}V ( M_{\,\text{возд}} - M_{\,\text{вод}} )} \cdot 100\%,\]

 то есть:

\[\varphi = \frac{160 \cdot 10 ^3 \cdot 30 \cdot 10 ^{ - 3 } \cdot 29 \cdot 10 ^{ - 3 } - 45 \cdot 10 ^{ - 3 } \cdot 8{,}31 \cdot 353}{47 \cdot 10 ^3 \cdot 30 \cdot 10 ^{ - 3 } \cdot ( 29 \cdot 10 ^{ - 3 } - 18 \cdot 10 ^{ - 3 } )} \cdot 100\%\approx 46\%.\]

Ответ: $46$\%.
}

% Source 43027: https://phys-ege.sdamgia.ru/problem?id=43027
\BankTask{43027}{
Общая масса газонепроницаемой оболочки воздушного шара и его корзины равна $400\,\text{кг}$. Шар заполнен гелием массой $100\,\text{кг}$. Он может поднять в воздух груз массой $225\,\text{кг}$. Какую минимальную массу гелия нужно добавить в оболочку шара, чтобы шар поднял еще одного пассажира массой $50\,\text{кг}$? Считать, что оболочка шара не оказывает сопротивления изменению его объема, воздушных течений в вертикальном направлении нет. Температура и давление гелия внутри шара и воздуха снаружи шара соответственно одинаковы.
}{
\smallskip\noindent\textbf{Решение}\par
Обозначим массу оболочки, гелия и груза $M=m_o + m_{\,\text{г}} + m.$ На шар действуют сила тяжести и выталкивающая сила со стороны воздуха. При поднятии шара подъемная сила равна 

\[F_{\,\text{под}}=F_A - Mg,\]

где выталкивающая сила $F_A=\rho_{\,\text{воз}}gV,$ V - объем шара. Плотность воздуха найдем из уравнения Менделеева-Клапейрона

$p= \frac{\rho_{\,\text{воз}}RT}{M_{\,\text{воз}}} ,$ откуда $\rho_{\,\text{воз}}= \frac{pM_{\,\text{воз}}}{RT} ,$

где p - давление наружного воздуха, T − температура воздуха снаружи шара.
Объем шара найдем из уравнения Менделеева-Клапейрона для гелия:

$pV= \frac{mRT}{M_{\,\text{г}}} ,$ откуда $V= \frac{mRT}{pM_{\,\text{г}}} ,$

где p - давление гелия, равное внешнему давлению, T - температура гелия, равная наружной температуре по условию.
Таким образом, подъемная сила шара равна:

\[F_{\,\text{под}}=g ( \frac{M_{\,\text{воз}}m}{M_{\,\text{г}}} - m_o - m - m_{\,\text{г}} ) =10 \cdot ( \frac{29 \cdot 10 ^{ - 3 } \cdot 100}{4 \cdot 10 ^{ - 3 }} - 400 - 100 - 225 ) =0.\]

Это означает, что масса поднимаемого груза является максимальной.
При добавлении пассажира и минимальной массы гелия в шар подъемная сила не изменится и будет выполняться аналогичное первому случаю соотношение:

\[m_o + m + m_{\,\text{г}} + m_{\,\text{п}}= \frac{M_{\,\text{воз}} ( m + \Delta m )}{M_{\,\text{г}}} .\]

При объединении соотношений для двух случаев получаем:

\[m_{\,\text{п}} + \Delta m= \frac{M_{\,\text{воз}}\Delta m}{M_{\,\text{г}}} ,\]

откуда находим минимальную массу гелия, который нужно добавить в шар:

\[\Delta m= \frac{m_{\,\text{п}}M_{\,\text{г}}}{M_{\,\text{воз}} - M_{\,\text{г}}} = \frac{50 \cdot 4 \cdot 10 ^{ - 3 }}{( 29 - 4 ) \cdot 10 ^{ - 3 }} =8\,\text{кг}.\]

Ответ: $8\,\text{кг}$.
}

% Source 33071: https://phys-ege.sdamgia.ru/problem?id=33071
\BankTask{33071}{
\begin{center}\begin{tikzpicture}[task graph,>=stealth]\draw(0,.8)--(0,2.7);\foreach \y in {.85,1.1,...,2.6}{\draw(0,\y)--++(-.22,-.21);}\draw(.21,2.09)circle(.19);\draw(.21,2.09)circle(.07);\draw(.21,2.01)rectangle(5,2.17);\node[above]at(.28,2.24){\(A\)};\node[above]at(1.86,2.2){\(B\)};\draw(1.77,2.01)rectangle(1.93,1.25);\draw(1.45,1.03)rectangle(2.23,1.25);\draw[line width=1.4pt](.53,0)--(.53,1.03)--(3.04,1.03)--(3.04,0);\draw(5,2.01)--(5,1.42);\draw[fill=orange!30](4.76,.86)rectangle(5.24,1.42);\end{tikzpicture}\end{center}
В цилиндр объемом $0{,}5\,\text{м}^{3}$ насосом закачивается воздух со скоростью $0{,}002\,\text{кг}$/c. В верхнем торце цилиндра есть отверстие, закрытое предохранительным клапаном. Клапан удерживается в закрытом состоянии стержнем, который может свободно поворачиваться вокруг оси в точке А (см. рис.). К свободному концу стержня подвешен груз массой $2\,\text{кг}$. Клапан открывается через $580\,\text{с}$ работы насоса, если в начальный момент времени давление воздуха в цилиндре было равно атмосферному. Площадь закрытого клапаном отверстия $5 \cdot 10 ^{ - 4 } \,\text{м} ^2 ,$ расстояние АB равно $0{,}1\,\text{м}$. Температура воздуха в цилиндре и снаружи не меняется и равна $300$ K. Определите длину стержня, если его считать невесомым.
}{
\smallskip\noindent\textbf{Решение}\par
Клапан откроется, когда избыточная сила F давления воздуха на клапан изнутри цилиндра сравняется с силой давления стержня на этот клапан. Если превышение давления воздуха в цилиндре над атмосферным $\Delta p,$ а площадь клапана S, то:

\[F=S \cdot \Delta p.\]

Сила действия стержня на клапан равна $mg \cdot \frac{L}{l} ,$ где m, L и l соответственно масса груза, длина стержня и длина его участка АВ. Итак, должно выполняться условие:

\[S \cdot \Delta p > \,\text{или} \,\text{равно} mg \cdot \frac{L}{l} .\]

Дополнительное давление воздуха определяется увеличением массы $\Delta m_{\,\text{в}}$ воздуха в цилиндре. Согласно уравнению Клапейрона  — Менделеева:

\[\Delta p= \frac{\Delta m_{\,\text{в}}}{MV} RT,\]

где M  — молярная масса воздуха. Поэтому условие открытия клапана имеет вид:

\[\frac{S\Delta m_{\,\text{в}}}{MV} RT > \,\text{или} \,\text{равно} mg \cdot \frac{L}{l} ,\]

 или 

\[l > \,\text{или} \,\text{равно} \frac{mgMVL}{SRT\Delta m_{\,\text{в}}} .\]

Если насос закачивает каждую секунду $\omega$ кг воздуха, то за время t он закачает в цилиндр массу $\Delta m_{\,\text{в}}=\omega t.$ Следовательно, клапан откроется в момент, когда выполнится равенство, из которого найдем длину стержня:

\[L= \frac{SRT\omega tl}{mgMV} = \frac{5 \cdot 10 ^{ - 4 } \cdot 8{,}31 \cdot 300 \cdot 0{,}002 \cdot 580 \cdot 0{,}1}{2 \cdot 10 \cdot 29 \cdot 10 ^{ - 3 } \cdot 0{,}5} \approx 0{,}5\,\text{м} .\]

Ответ: $0{,}5\,\text{м}$.
}

% Source 29749: https://phys-ege.sdamgia.ru/problem?id=29749
\BankTask{29749}{
Закрытый вертикальный цилиндрический сосуд, заполненный идеальным газом, разделён на две части тяжёлым поршнем, способным скользить без трения. В начальном равновесном состоянии в верхней и нижней частях сосуда находилось по $\nu$  =  $1$ моль газа, а отношение объёмов верхней и нижней частей сосуда было равно $2$. После того как из верхней части сосуда полностью откачали газ, через длительный промежуток времени установилось новое состояние равновесия. Найдите отношение объёмов верхней и нижней частей сосуда после откачки газа. Температура газа T в обеих частях сосуда всё время поддерживалась одинаковой и постоянной.
}{
\smallskip\noindent\textbf{Решение}\par
$1$.  Вначале на поршень действовали: сила давления газа в верхней части цилиндра $F_1=p_1 S,$ направленная вниз, сила давления газа в нижней части цилиндра $F_2=p_2 S,$ направленная вверх, сила тяжести поршня mg, направленная вниз, где $p_{1}$  и $p_{2}$  — давления газа в верхней и нижней частях цилиндра, S  — площадь поршня, m  — масса поршня. Поскольку поршень находится в равновесии, то $p_2 S=p_1 S + m g.$ Учитывая, что давление в верхней части меньше, чем в нижней, делаем вывод, что объем газа в верхней части больше, чем объем газа в нижней части цилиндра. 
Обозначим объем нижней части цилиндра V, объем верхней части цилиндра 2V.
Используя уравнение Клапейрона  — Менделеева $p V= \nu R T,$ выразим давления каждой части газа: $p_1= \frac{\nu R T}{2 V}$ и $p_2= \frac{\nu RT}{V} .$ Тогда сила тяжести, действующая на поршень $m g= \frac{\nu RTS}{2 V} .$ 
$2$.  Когда из верхней части цилиндра удалили весь газ, поршень поднимался до тех пор, пока сила тяжести не уравновесилась силой давления газа, находящегося в нижней части цилиндра, то есть $m g=p S.$ С учетом выраженного значения силы тяжести и применяя уравнение Клапейрона  — Менделеева для этого случая, получаем объем газа в нижней части цилиндра:

\[V_2= \frac{\nu R T}{p} = \frac{\nu R T \cdot 2 V}{\nu R T} =2 V.\]

Таким образом, в нижней части объем газа будет в $2$ раза больше, чем в верхней. Значит, отношение объема верхней части цилиндра к объему нижней части станет равным $0{,}5$. 

Ответ: $0{,}5$.
}

% Source 9248: https://phys-ege.sdamgia.ru/problem?id=9248
\BankTask{9248}{
Мальчик решил подняться в воздух на воздушных шарах с гелием. Известно, что мальчик весит $40\,\text{кг}$, а на улице нормальное атмосферное давление и температура $27$ °C. Учитывая, что объем одного шара составляет $10$ литров, найдите, сколько потребуется шаров для такого путешествия. Массой оболочки шаров и объемом мальчика пренебречь.
}{
\smallskip\noindent\textbf{Решение}\par
На мальчика действуют силы: сила Архимеда $F_A=\rho_{\mathrm{He}} ( nV ) g,$ сила тяжести, действующая на человека, mg и сила тяжести, действующая на шарики, $m_{\mathrm{He}}g=\rho_{\mathrm{He}}g ( nV ) .$
Мальчик будет подниматься за счет силы Архимеда, действующей на шары:

$mg + m_{\mathrm{He}}g=F_A.$
$mg + \rho_{\mathrm{He}} ( nV ) g=g\rho_{\,\text{в}} ( nV ) ,$

откуда количество шариков:

\[n= \frac{m}{V ( \rho_{\,\text{в}} - \rho_{\mathrm{He}} )} ,\]

где $\rho_{\,\text{в}}$  — плотность воздуха, $\rho_{\mathrm{He}}$  — плотность гелия.
Из уравнения Менделеева  — Клапейрона можем найти плотность газа при данной температуре и давлении:

$p= \frac{\rho}{M} RT ,$
$\rho= \frac{pM}{RT} .$

Найдем необходимое количество шаров, учитывая, что гелий в шариках находится при той же температуре и давлении, что и воздух вокруг:

\[n= \frac{m}{V ( \rho_{\,\text{в}} - \rho_{\mathrm{He}} )} = \frac{mRT}{Vp_0 ( M_{\,\text{в}} - M_{\mathrm{He}} )} = \frac{40 \cdot 8{,}31 \cdot 300}{10 \cdot 10 ^{ - 3 } \cdot 10 ^{ 5 } \cdot ( 0{,}029 - 0{,}004 )} \approx 4000.\]

Ответ: $4000$.
}

% Source 7165: https://phys-ege.sdamgia.ru/problem?id=7165
\BankTask{7165}{
\begin{center}\begin{tikzpicture}[task graph,>=stealth]\draw(0,1.45)circle(1.45);\draw(-1.3,.82)--(-.48,-1.18)--(.48,-1.18)--(1.3,.82);\draw(-.58,-.91)rectangle(.58,-1.18);\foreach \x in {-.5,-.34,...,.46}{\draw(\x,-.91)--++(.06,-.27);}\draw(-.09,-.91)--(-.09,-.14)..controls(-.7,.15)and(-.4,.34)..(-.56,.57)..controls(-.71,.91)and(-.18,.97)..(-.23,1.2)..controls(-.3,1.43)and(.27,1.64)..(.3,1.28)..controls(.57,1.11)and(.38,.83)..(.51,.65)..controls(.74,.35)and(.32,.3)..(.09,-.12)--(.09,-.91);\node at(.74,.9){\(t\)};\node at(-1.31,-.26){\(t_0\)};\end{tikzpicture}\end{center}
Воздушный шар, оболочка которого имеет массу M  =  $145\,\text{кг}$ и объем V  =  $230\,\text{м}^{3}$, наполняется горячим воздухом при нормальном атмосферном давлении и температуре окружающего воздуха $t_{0}$  =  $0$ °С. Какую минимальную температуру t должен иметь воздух внутри оболочки, чтобы шар начал подниматься? Оболочка шара нерастяжима и имеет в нижней части небольшое отверстие.
}{
\smallskip\noindent\textbf{Решение}\par
На шар действуют: сила Архимеда $F_A=\rho_{\,\text{в}}gV,$ сила тяжести, действующая на воздух внутри, $mg$ и сила тяжести, действующая на оболочку, $Mg.$ По второму закону Ньютона:

\[\vec F_A + m \vec{g} + M \vec{g}=0.\]

В  проекциях на вертикаль: $F_A=mg + Mg.$
Условие подъема шара: $F_{\,\text{Арх}}. > \,\text{или} \,\text{равно} Mg + mg,$ где M  — масса оболочки, $m=\rho V$— масса воздуха внутри оболочки, отсюда:

$\rho_0gV > \,\text{или} \,\text{равно} Mg + \rho gV,$
$\rho_0V > \,\text{или} \,\text{равно} M + \rho V,$

где \ensuremath{\rho}$_{0}$  — плотность окружающего воздуха, \ensuremath{\rho} — плотность воздуха внутри оболочки, V  — объем шара.
Для воздуха внутри шара находим: $\frac{pV}{T} = \frac{m}{\mu} R,$ или $\frac{m}{V} = \frac{p \cdot \mu}{R \cdot T} = \rho,$ где p  — атмосферное давление, $\mu$  — молярная масса воздуха, T  — температура воздуха внутри шара. Соответственно, имеем плотность воздуха снаружи: $\rho_0= \frac{\mu p}{RT_0} ,$ где $T_{0}$  — температура окружающего воздуха.

$\frac{p \cdot \mu \cdot V}{R \cdot T_0} > \,\text{или} \,\text{равно} M + \frac{p \cdot \mu \cdot V}{R \cdot T} .$
$\frac{p \cdot \mu \cdot V}{R \cdot T_{\min}} = \frac{p \cdot \mu \cdot V}{R \cdot T_0} - M.$
$\frac{1}{T_{\min}} = \frac{1}{T_0} - \frac{M \cdot R}{p \cdot \mu \cdot V} ,$

$T_{\min}=T_0 \frac{p\mu V}{p\mu V - MRT_0} =273 \cdot \frac{10 ^{5} \cdot 29 \cdot 10 ^{ - 3 } \cdot 230}{10 ^{5} \cdot 29 \cdot 10 ^{ - 3 } \cdot 230 - 145 \cdot 8{,}31 \cdot 273} \approx 538\,\text{К} = 265 ^\circ \,\text{С}.$

Ответ: $265$.
}

% Source 7935: https://phys-ege.sdamgia.ru/problem?id=7935
\BankTask{7935}{
Два одинаковых теплоизолированных сосуда соединены короткой трубкой с краном. Объем каждого сосуда V  =  $1\,\text{м}^{3}$. В первом сосуде находится $\nu_{1}$  =  $1$ моль гелия при температуре $T_{1}$  =  $400$ К; во втором  — $\nu_{2}$  =  $3$ моль аргона при температуре $T_{2}$. Кран открывают. После установления равновесного состояния давление в сосудах р  =  $5{,}4$ кПа. Определите первоначальную температуру аргона $T_{2}$.
}{
\smallskip\noindent\textbf{Решение}\par
$1$.  Поскольку в указанном процессе газ не совершает работы и система является теплоизолированной, то в соответствии с первым законом термодинамики суммарная внутренняя энергия газов сохраняется:

\[\frac{3}{2} \nu_1RT_1 + \frac{3}{2} \nu_2RT_2= \frac{3}{2} ( \nu_1 + \nu_2 ) RT,\]

где T  — температура в объединенном сосуде в равновесном состоянии после открытия крана.
$2$.  В соответствии с уравнением Клапейрона  — Менделеева для конечного состояния можно записать: 

\[p ( 2V ) = ( \nu_1 + \nu_2 ) RT.\]

Исключая из двух записанных уравнений конечную температуру T, получаем искомое выражение для начальной температуры аргона:

\[T_2= \frac{2Vp}{\nu_2R} - \frac{\nu_1}{\nu_2} T_1= \frac{2 \cdot 1 \cdot 5{,}4 \cdot 10 ^3}{3 \cdot 8{,}31} - \frac{1}{3} \cdot 400 \approx 300\,\text{К}.\]

Ответ: $T_2 \approx 300\,\text{К}.$
}

% Source 9976: https://phys-ege.sdamgia.ru/problem?id=9976
\BankTask{9976}{
Объем $0{,}1$ литра водорода нагревают при постоянном давлении от $300$ до $3000$ К. При высоких температурах молекулы водорода распадаются на отдельные атомы. На графике показана зависимость доли распавшихся молекул от температуры. Чему равен конечный объем газа?
\begin{center}\begin{tikzpicture}[task graph,>=stealth]\draw[gray!65,xstep=.6,ystep=.6](0,0)grid(5.4,3.6);\draw[->](0,0)--(5.4,0)node[below]{\(T,\,\text{К}\)};\draw[->](0,0)--(0,3.7)node[left]{\(\alpha,\,\%\)};\node[below left]at(0,0){\(0\)};\foreach \x/\t in {1.2/1000,2.4/2000,3.6/3000}{\node[below]at(\x,0){\(\t\)};}\foreach \y/\a in {.6/10,1.2/20,1.8/30,2.4/40}{\node[left]at(0,\y){\(\a\)};}\draw[smooth] plot coordinates{(0,0)(.6,.008)(1.2,.052)(1.8,.16)(2.4,.34)(3,.69)(3.6,1.2)(4.2,1.81)(4.8,2.53)(5.4,3.3)};\end{tikzpicture}\end{center}
}{
\smallskip\noindent\textbf{Решение}\par
В первоначальном состоянии весь водород находился в молекулярном состоянии, и согласно уравнению Менделеева  — Клапейрона:

\[pV_1=\nu RT_1.\]

После нагревания согласно графику $\alpha =20\%$ молекул распадется. При распаде из одной молекулы водорода образуется два атома водорода, значит, в сосуде будет $( 1 - \alpha ) \nu$ молекулярного водорода и $2 \alpha \nu$ атомарного водорода. Общее количество вещества стане равным $( 1 - \alpha ) \nu + 2 \alpha \nu = ( 1 + \alpha ) \nu = 1{,}2\nu,$ и согласно уравнению Менделеева  — Клапейрона:

\[pV_2=1{,}2\nu RT_2.\]

Поделив второе уравнение на первое, получаем:

$\frac{V_2}{V_1} =1{,}2 \frac{T_2}{T_1} =1{,}2 \cdot \frac{3000}{300} =12,$

$V_2=12V_1=12 \cdot 0{,}1=1{,}2\,\text{л}.$

Ответ: $V_2=1{,}2\,\text{л}.$
}

\end{document}
